\section{Extrema af funktioner}

\subsection{Generelt}

Hvis en funktion $f$ er:
\begin{itemize}
    \item Kontinuert
    \item Defineret på et lukket og begrænset interval $[a,b]$
\end{itemize}
Så er $\image(f)$ også begrænset, og har både minimal- og maksimalværdi.

\[
f([a,b]) = [min, Max]
\kildetag{Theorem 5.1.1}
\]

Desuden vil der for hver $y \in [min, Max]$ være mindst ét $x \in [a,b]$ så $f(x) = y$.

\subsection{Supremum og infimum}
Et Supremum for en mængde $A$ er det mindste tal $s$ så $\forall a \in A : a \geq a$.

Et Infimum for en mængde $A$ er det største tal $i$ så $\forall a \in A : a \leq i$.

\eksempel{
    For en mænde $A = ]6,7[$ er $\sup(A) = 7$ og $\inf(A) = 6$. 
    
    Bemærk at hverken $6$ eller $7$ er elementer i $A$, og derfor ikke er maksimum eller minimum for $A$.
    ($A$ har ingen maksimum eller minimum).
}

\subsection{Hvor finder man ekstrema?}
De eneste steder et ekstrema i intervallet $[a,b]$ kan være er:
\begin{enumerate}
    \item $a$ (venstre endepunkt)
    \item $b$ (højre endepunkt)
    \item Punkter hvor $f'(x) = 0$ (stationære punkter)
    \item Punkter hvor $f'(x)$ ikke er defineret (singulære punkter)
\end{enumerate}

\eksempel{
    Find ekstrema for $f: [0,2] \to \mathbb R, f(x) = x^3 - x$
    
    Først de simple:
    \[
        f(0) = 0, f(2) = 6
    \]

    Funktionen er kontinuert, så vi kan ignorere tilfælde 4.

    For at finde stationære punkter differentierer vi $f$:
    \(
        f'(x) = 3x^2 - 1
    \)

    Vi skal nu finde punkt(erne) hvor $f'(x) = 0$:
    \[
        3x^2 - 1 = 0 \implies 3x^2 = 1 \implies x^2 = \frac{1}{3} \implies x = \pm \frac{1}{\sqrt{3}}
    \]

    Kun $x = +\frac{1}{\sqrt{3}}$ er i intervallet $[0,2]$.
    Vi har nu fundet alle kandidater og evaluerer $f$ i dem:
    \[
        f(0) = 0, f(2) = 6, f\left(\frac{1}{\sqrt{3}}\right) = \frac{1}{3\sqrt{3}} - \frac{1}{\sqrt{3}} = -\frac{2}{3\sqrt{3}}
    \]

}

\subsection{Vurdering af stationære punkter}
For at vurdere et stationært punkt, tag anden afledte $f''(x)$:
\begin{itemize}
    \item Hvis $f''(x) > 0 \implies$ lokalt minimum
    \item Hvis $f''(x) < 0 \implies$ lokalt maksimum
    \item Hvis $f''(x) = 0$ usikkert\footnote{Her kan du så fortsætte med at tage tredje afledte, og så videre, indtil du finder en afledte der ikke er 0. Hvis den første afledte der ikke er 0 er af lige orden, så er det et lokalt minimum hvis den er positiv, og et lokalt maksimum hvis den er negativ. Hvis den første afledte der ikke er 0 er af ulige orden, så er det hverken et lokalt minimum eller maksimum.}
\end{itemize}

\kilde{Theorem 5.1.3}

\subsection{Ekstrema i flere variable}
Hvis en funktion $f$ er defineret på et lukket og begrænset område $A \subseteq \mathbb R^n$ og er kontinuert,\footnote{For ethvert valg af 2 punkter i $A$, findes en kontinuert kurve imellem der ligger i $A$.}
så har $f$ både minimum og maksimum i $A$.

Ekstrema kan kun findes i:
\begin{itemize}
    \item $A \cap \partial A$ (grænsen af $A$)
    \item $A^\circ$ hvor $f$ ikke er differentiabel
    \item $A^\circ$ hvor $\nabla f = \vec 0$
\end{itemize}

\kilde{Theorem 5.2.2}

\eksempel{
    Find statinære punkter for $f: \mathbb R^2 \to \mathbb R, f(x,y) = xy^2 + x^2 + y^3$.

    Find først gradienten $\nabla f$:
    \[
        \nabla f(x,y) = 
        \begin{bmatrix}
            y^2 + 2x &
            2xy + 3y^2
        \end{bmatrix}^T
    \]

    Fra første indgang kræves det at $2x = -y^2$ for at det bliver 0. Indsætter vi det i anden indgang, får vi:
    \[
        -y^3 + 3y^2 = 0 \implies y^2(-y + 3) = 0 \implies y = 0 \lor y = 3
    \]

    Sætter vi de to ind i første indgang får vi:
    \[
        0^2+2x = 0 \implies x = 0 \qquad
        3^2 + 2x = 0 \implies x = -\frac{9}{2}
    \]

    Derfor er de stationære punkter $(0,0)$ og $(-\frac{9}{2}, 3)$.
}

\eksempel{
    Find ekstrema i $\partial A$.
    \[
    f: \mathbb B \to \mathbb R, \quad
    B \in \{ (x,y) \in \mathbb R^2 | x^2+y^2 \le 4 \}, \quad
    f(x,y) = xy + 1
    \]

    Randen består af cirklen med radius 2, og kan derfor parametriseres som:
    $\vec r(t) = [2\cos(t), 2\sin(t)]^T, t \in [0, 2\pi]$

    Finder vi nu $f \circ \vec r$ får vi:
    \[
        (f \circ \vec r)(t) = 2 \cdot 2 \cdot \cos(t) \cdot \sin(t) + 1 = 2 \sin(2t) + 1
    \]

    $\sin(2t), t \in [0, 2\pi]$ har intervallet $[-1, 1]$, derfor har randen ekstrema: $-1, 3$
}

\subsection{Lokale ekstrema}

Givet rum $A \subseteq \mathbb R^n$ og funktion $f: A \to \mathbb R$:

\begin{align*}
    \exists \epsilon > 0 : \forall \vec x \in A \setminus \{\vec x_0\}: ||\vec x - \vec x_0|| < \epsilon &\implies f(\vec x) \le f(\vec x_0)
    \qquad \text{(lokalt maksimum)} \\
    \exists \epsilon > 0 : \forall \vec x \in A \setminus \{\vec x_0\}: ||\vec x - \vec x_0|| < \epsilon &\implies f(\vec x) \ge f(\vec x_0)
    \qquad \text{(lokalt minimum)}
\end{align*}

Hvis man kan udskifte $\le$ med $<$, og $\ge$ med $>$, så er de \emph{strenge} ekstrema. \\

Alle lokale ekstrema er stationære punkter ($\nabla f = \vec 0$) eller punkter hvor $f$ ikke er differentiabel.

\kilde{Lemma 5.2.3}

\subsection{Find lokale ekstrema}

Givet skalar funktion $f: U \to \mathbb R$, hvor $U \subseteq \mathbb R^n$ er en åben mængde, og $f$ er 2-gange differentiabel i $\vec x_0 \in U$:

Lav en hesse-matrix $H_f(\vec x_0)$ (Se \autoref{sec:HesseMatricer}) og vurder den:
\begin{itemize}
    \item $H_f(\vec x_0)$ er positivt definit $\implies$ strengt lokalt minimum
    \item $H_f(\vec x_0)$ er negativt definit $\implies$ strengt lokalt maksimum
    \item $H_f(\vec x_0)$ er indefinit $\implies$ saddelpunkt
    \item $H_f(\vec x_0)$ er semidefinit $\implies$ usikkert
\end{itemize}
\kilde{Theorem 5.2.4}